\chapter{Controle climático e calibração da referência}
\label{cap:calibracao}

O capítulo anterior mede uma queda de fator de capacidade sem separar restrição
de variação do recurso. A irradiância medida sobre cada usina faz essa
separação, e a separação depende de um pressuposto sobre o ano de base.

\section{Restrição ou recurso}
\label{sec:fig5}

\begin{ficha}
  \item[Janela]      2023 a julho de 2026
  \item[Amostra]     meses com cobertura climática
  \item[Unidade]     conjunto por mês
  \item[Fonte]       ONS; NASA POWER, irradiância diária no ponto da usina
  \item[Teste]       nenhum: decomposição aritmética
\end{ficha}

Com a irradiância global horizontal diária $H$ sobre a usina, a razão de
desempenho divide a energia gerada por kW instalado pelo recurso disponível:

\[ PR = \frac{FC \cdot 24\,\mathrm{h}}{H} \]

Se a queda do fator de capacidade viesse de menor irradiância, $PR$ permaneceria
estável.

\begin{figure*}[t]
\centering
\includegraphics[width=\linewidth]{../figures/fig5_controle_climatico.pdf}
\caption{\textbf{Normalizada pela irradiância medida sobre a usina, a perda de
Sol do Cerrado permanece; o recurso solar não a explica.}
\textbf{a}, Razão de desempenho mensal, definida como energia gerada por kW
instalado dividida pela irradiância global horizontal diária. O tracejado
horizontal marca a média de cada ano.
\textbf{b}, Irradiância global horizontal mensal sobre cada usina, em linha
cheia, e a mesma sob céu claro, em tracejado.
\textbf{c}, Decomposição da queda do fator de capacidade de Sol do Cerrado entre
2023 e 2025: observado em 2023, esperado em 2025 caso apenas a irradiância
tivesse mudado, e observado em 2025.
\textbf{d}, Fator de capacidade mensal contra irradiância mensal em Sol do
Cerrado, por ano.}
\label{fig:clima}
\end{figure*}

A irradiância sobre Jaíba caiu 2,8\% entre 2023 e 2025, e o fator de capacidade
caiu 29,1\%. Escalando o fator de 2023 pela razão de irradiância, o esperado
para 2025 seria 0,206, e o observado foi 0,150. Cerca de um décimo da queda é
atribuível ao recurso. A razão de desempenho cai de 0,854 para 0,619 em Sol do
Cerrado, enquanto o controle Três Marias 3 permanece em torno de 0,99 sem
tendência.

\conclusao{Nove décimos da queda do fator de capacidade não são atribuíveis à
irradiância. A perda é de rede, e não de recurso.}

Os meses de 2025 não estão sobre a curva de 2023 em posição de menor
irradiância. Eles estão sobre uma curva mais baixa. O deslocamento é consistente
com perda independente do recurso.

\campo{Uma conclusão secundária que não se sustentou} O déficit não explicado
pela irradiância é de 0,056 em fator de capacidade, e o corte apurado pelo ONS
equivale a 0,147, ou 2,6 vezes mais. A diferença foi interpretada como indício
de que a geração de referência do operador superestima o contrafactual. A
interpretação pressupõe que 2023 fosse um ano sem restrição. A
seção~\ref{sec:fig6} mostra que não era.

\begin{objecao}
A irradiância vem de reanálise, e não de piranômetro: NASA POWER combina MERRA-2
com satélite em célula de 0,5 por 0,625 grau. A grandeza é irradiância
horizontal, e os módulos não são horizontais. A razão de desempenho é indicador
relativo dentro de cada usina, e não a razão de norma IEC. A
temperatura não é modelada e o contrafactual é linear.
\end{objecao}

\campo{Erro de variável derivada de sensoriamento} A irradiância entra como
denominador. Proctor, Carleton e Sum \cite{proctor2023} mostram que uma variável
assim enviesa o ponto e o erro padrão de uma inferência a jusante, mesmo com
erro de medida pequeno. Eles recomendam tratá-la como distribuição, por
imputação múltipla. Alix-García e Millimet \cite{alixgarcia2023} acrescentam que o erro de
produto de satélite raramente é clássico. Nenhuma das duas correções foi
aplicada. A diferença medida é grande o bastante para sobreviver a erro
moderado, e o intervalo em torno dela é mais estreito do que deveria ser.

\section{Calibração da geração de referência}
\label{sec:fig6}

\begin{ficha}
  \item[Janela]      abril de 2024 a julho de 2026
  \item[Amostra]     74 usinas com 30 dias limpos ou mais
  \item[Unidade]     usina; dia-usina na calibração
  \item[Fonte]       ONS; NASA POWER
  \item[Teste]       Spearman sobre 2 covariáveis
\end{ficha}

Para cada usina calibra-se uma razão de desempenho nos dias sem restrição da
própria usina, com fração cortada abaixo de 1\%. Compara-se então o potencial
implicado pelo ONS, igual à geração verificada mais a cortada, com o potencial
previsto pela irradiância local sob aquela razão calibrada.

\begin{figure*}[t]
\centering
\includegraphics[width=\linewidth]{../figures/fig6_vies_referencia.pdf}
\caption{\textbf{A geração de referência do ONS é bem calibrada; o viés aparente
da seção anterior vinha do contrafactual, e não do operador.}
\textbf{a}, Distribuição do viés mediano por usina, definido como a razão entre
o potencial implicado pelo ONS e o potencial previsto pela irradiância local sob
a razão de desempenho calibrada nos dias sem restrição da própria usina. O
tracejado em 1 marca ausência de viés.
\textbf{b}, Viés contra a intensidade mediana de corte nos dias restritos.
\textbf{c}, Validação da calibração: fator de capacidade contra irradiância nos
dias sem restrição de Sol do Cerrado, com a reta ajustada pela razão de
desempenho limpa.
\textbf{d}, Razão de desempenho anual de Sol do Cerrado contra o desempenho
limpo calibrado. O rótulo interno indica o corte implícito de cada ano.}
\label{fig:calibracao}
\end{figure*}

O viés mediano por usina é de 1,02, com metade das usinas entre 0,98 e 1,06.
Nenhuma das 74 usinas calibradas passa de 1,5, e 26 ficam abaixo de 1, o que
torna a referência conservadora para um terço delas. No agregado energético dos
dias restritos, o corte apurado é 1,03 vezes o implicado pela irradiância. A
associação entre viés e intensidade de corte é nula, o que descarta a falha de
maior consequência.

\conclusao{A geração de referência do ONS não apresenta viés detectável, e não
se degrada onde o corte é mais intenso.}

\campo{A revisão} O desempenho limpo de Sol do Cerrado, calibrado em 279 dias
sem restrição, é de 0,931. Em 2023 a usina operou a 0,854, ou 8,3\% abaixo do
seu próprio desempenho limpo. Aquele ano já era restrito, e não havia apuração
pública. Ao usar 2023 como base, a seção~\ref{sec:fig5} mediu um déficit
reduzido e concluiu de forma incorreta que o operador superestimava o corte.

\campo{O que muda e o que permanece} A conclusão principal da
seção~\ref{sec:fig5} permanece e sai reforçada. Contra o desempenho limpo, o
corte implícito é de 33\% em 2025, e não de 26\%. A afirmação secundária sobre
viés do operador é retirada. A taxa agregada da seção~\ref{sec:fig1} não requer
correção, porque uma referência calibrada sustenta o agregado.

\begin{objecao}
Dias classificados como limpos podem não ser limpos. Como o resultado é ausência
de viés, esse erro só o tornaria mais forte. Seis usinas ficaram de fora por não
atingirem 30 dias limpos, e são as mais restritas, onde a referência é mais
usada. A razão de desempenho limpa é constante por usina e herda qualquer viés
sazonal dos dias de calibração.
\end{objecao}

\campo{Posição na literatura} A geração de referência do ONS é derivada de
curvas de desempenho da própria usina, como descrevem Vieira et al.
\cite{vieira2026}. Validá-la exige fonte independente. No lado eólico, Silva et
al. \cite{silva2026} usam anemometria de SCADA. A verificação equivalente no
lado solar, com irradiância de satélite, não foi localizada na literatura
consultada, e a busca não foi sistemática.
